% HD-G15. 21 November 2026. GPT-6 Astra Pro, inherited effort.
\section{The fourfold point ideal and its explicit Fourier square}

\paragraph{HD-G15.}
The ideal of one point on an abelian fourfold is an actual
indecomposable perfect object whose trace is injective in
every degree. Its external square has degree-two self-Ext
dimension88, so the rank calculation HD-N14 alone does not
exclude it. The explicit equal-twist Fourier square described
below nevertheless cannot have whole character
\[
                  e^B\bigl(Q(\theta)+\lambda t\bigr)       \tag{G15.1}
\]
for any compatible ample Weil-type polarization $\theta$,
any nonzero Weil class $t$ on the eightfold, any polynomial
$Q$, and any algebraic rational $B$. The proof computes the
whole character of this actual object and exhibits a
nonzero component of its centered second character.

\paragraph{The actual point-ideal cone and all its traces.}
Let $X$ be an abelian fourfold, $p\in X$, and $I=I_p$.
The exact triangle
\[
                  I\longrightarrow\mathcal O_X
                      \longrightarrow\mathcal O_p
\]
constructs this object as an actual cone, shifted by $-1$.
It is a sheaf, so it has no negative self-Ext. It has rank
one and double dual $\mathcal O_X$. Every endomorphism
extends to the double dual and is scalar; thus it is simple
and indecomposable.

The local codimension-four Koszul resolution and the global
long exact sequence give
\[
 0\longrightarrow T_pX\longrightarrow\operatorname{Ext}^1(I,I)
        \longrightarrow H^1(X,\mathcal O_X)\longrightarrow0.
                                                               \tag{G15.2}
\]
Here $\operatorname{Hom}(I,\mathcal O_p)=T_pX$,
$\operatorname{Ext}^1(I,\mathcal O_X)=H^1(\mathcal O_X)$,
and the latter group's map to
$\operatorname{Ext}^1(I,\mathcal O_p)$ is zero. Indeed it
comes from scalar cohomology of $\mathcal O_X$; centrality
of that action and the zero composite
$I\to\mathcal O_X\to\mathcal O_p$ give the zero map.
The map
$\operatorname{Hom}(I,I)\to\operatorname{Hom}(I,\mathcal O_X)$
is the isomorphism between the two scalar lines, so(G15.2)
has no omitted term.
Consequently $\operatorname{ext}^1(I,I)=8$.

Its character is $1-[p]$, and its derived-dual character
is also $1-[p]$, since the point has codimension four.
Riemann--Roch gives $\chi(I,I)=-2$. Serre duality now gives
the complete array
\[
       (\operatorname{ext}^0,\operatorname{ext}^1,
          \operatorname{ext}^2,\operatorname{ext}^3,
          \operatorname{ext}^4)(I,I)=(1,8,12,8,1).
                                                               \tag{G15.3}
\]
There are no further degrees.

For contraction with $1-[p]$, the four antiholomorphic
wedges detect the rank term and the four holomorphic
contractions detect the point term in a different Hodge
summand. Thus degree one has rank8. In degree two, the
six double wedges and six double contractions again
have independent images; all sixteen mixed operators
kill both terms. Thus degree two has rank12.
Comparing these ranks with(G15.3) proves injectivity of
the degree-one and degree-two traces and surjectivity
of their evaluation maps. The degree-three trace is
transpose to degree-one evaluation under fourfold
Serre duality, so it injects as well. The rank-one
identity character and its Serre-dual identity
evaluation give injectivity in degrees zero and four.
This proves the assertion about every trace degree.
Tensoring with an actual line bundle and shifting
preserve the assertion by the trace functoriality
and invertible Chern-character shear already proved.

For $I\boxtimes I$, the actual K\"unneth and tensor-trace
maps give injectivity in every degree and
\[
       \operatorname{ext}^1=16,\qquad
       \operatorname{ext}^2=12+8^2+12=88.                 \tag{G15.4}
\]
The degree-six/degree-two adjointness of HD-N14 therefore
gives contraction rank88, the allowed value with
$\rho=2$ in that result.

\paragraph{The specified Fourier construction.}
Choose the actual CM fourfold $X=E_i^4$, its origin $0$,
and the product principal line bundle with class
$H=\sum_{j=1}^4a_jb_j$. Set $L_\ell=\mathcal O_X(\ell H)$
for an integer $\ell$, and put
\[
 F_\ell=(I_0\otimes L_\ell)\boxtimes(I_0\otimes L_\ell).
 \qquad
 \mu(x,y)=(x+y,y).
\]
Use the normalized Poincare bundle $\mathcal P$ and the
actual enhanced equivalence
\[
 \Phi=(\operatorname{id}_X\boxtimes
                 \operatorname{FM}_{\mathcal P^{-1}[4]})\circ\mu^*,
 \qquad \mathcal E_\ell=\Phi(F_\ell).
                                                               \tag{G15.5}
\]
The Fourier equivalence is Mukai's Theorem2.2,
Nagoya Mathematical Journal81 (1981), p.156.
The output is $X\times\widehat X$, identified with a product
of eight CM elliptic curves using the product principal
polarization. Write
\[
 A=H_x,\quad D=H_{\widehat x},\quad
 P=c_1(\mathcal P_{x,\widehat x}),\quad
 \pi_x=[0]_x=H_x^4/4!,\quad
 \pi_{\widehat x}=H_{\widehat x}^4/4!.
\]
We compute the whole character
$v_\ell=e^{-P/2}\operatorname{ch}(\mathcal E_\ell)$.
This multiplication is an operation on rational cohomology;
no fractional tensor functor is asserted.

For $\ell\ne0$ the exact answer is
\[
 \boxed{\quad
 v_\ell=(2\ell)^4
          \exp\left(\frac{\ell}{2}A-\frac{1}{2\ell}D\right)
       -\exp(\ell A+P/2)-\exp(\ell A-P/2)+\pi_x .
       \quad}                                               \tag{G15.6}
\]
For $\ell=0$ it is
\[
           v_0=\pi_x+\pi_{\widehat x}
                         -\exp(P/2)-\exp(-P/2).             \tag{G15.7}
\]
All exponentials are truncated in the cohomology of the
eightfold.

Here is the explicit exterior integration proving both
formulas. For one elliptic coordinate use input one-forms
$a,b,c,d,u,v$ and set
\[
 h_{x+y}=(a+c)(b+d),\qquad h_y=cd,\qquad
 P_{\mathrm{in}}=cv-du,\qquad P_{\mathrm{out}}=av-bu .
\]
The coefficient of $cd$ is integration along the middle
elliptic curve. Direct exterior expansion gives
\[
\begin{aligned}
 \int_{c,d}\exp(\ell h_{x+y}+\ell h_y-P_{\mathrm{in}})
    &=2\ell\exp\left(\ell ab/2+P_{\mathrm{out}}/2
                                     -uv/(2\ell)\right),\\
 \int_{c,d}h_{x+y}\exp(\ell h_y-P_{\mathrm{in}})
    &=\exp(\ell ab+P_{\mathrm{out}}),\\
 \int_{c,d}h_y\exp(\ell h_{x+y}-P_{\mathrm{in}})
    &=\exp(\ell ab),\\
 \int_{c,d}h_{x+y}h_y&=ab .
\end{aligned}                                                \tag{G15.8}
\]
After multiplying the first expression by
$e^{-P_{\mathrm{out}}/2}$ it is the polynomial
\[
             2\ell+\ell^2ab-uv-\frac{\ell}{2}abuv.           \tag{G15.9}
\]
Thus its value at $\ell=0$ is $-uv$. Taking four products
of the four identities(G15.8), with signs $+,-,-,+$
from the two point ideals, proves(G15.6)--(G15.7);
the fourfold shift in(G15.5) has positive Chern-character
sign. This is an actual computation of the whole character,
including all terms arising from the two point factors.

\paragraph{The centered second character.}
For any character $v$ of nonzero rank $r$ set
\[
                    C(v)=v_2-\frac{v_1^2}{2r},              \tag{G15.10}
\]
where $v_j$ is its codimension-$j$ component.
Expanding(G15.6), including the polynomial specialization
at zero, gives
\[
\begin{aligned}
 r&=16\ell^4-2\ne0,\\
 (v_\ell)_1&=(8\ell^5-2\ell)A-8\ell^3D,\\
 (v_\ell)_2&=(2\ell^6-\ell^2)A^2-4\ell^4AD
                       +2\ell^2D^2-\tfrac14P^2.
\end{aligned}                                                \tag{G15.11}
\]
The rank is nonzero for every integer $\ell$.
The class $C(v_\ell)$ is a linear combination of the four
displayed quadratic classes. Its coefficient on $P^2$ is
$-1/4$. It is invariant under the actual automorphism
\[
       \iota=\operatorname{id}_X\times[-1]_{\widehat X}.
\]

Moreover its component in
\[
 H^{2,0}(X)\otimes H^{0,2}(\widehat X)
          \ \subset\ H^4(X\times\widehat X,\mathbb C)       \tag{G15.12}
\]
is nonzero. To see it directly, put
$z_j=a_j+ib_j$ and $w_j=u_j+iv_j$. Then
\[
       P=\frac{i}{2}\sum_j
             (z_j\bar w_j-\bar z_jw_j).
\]
For two distinct indices, the product of the
$z_j\bar w_j$ terms is a nonzero component of $P^2$
in(G15.12). None of $A^2,AD,D^2$ has that component.
For example, use the exterior order
\[
       z_1,\bar z_1,w_1,\bar w_1,z_2,\bar z_2,w_2,\bar w_2,\ldots.
\]
The coefficient of
$z_1\bar w_1z_2\bar w_2$ in $C(v_\ell)$ is exactly $1/8$.

\paragraph{The ample-square contradiction.}
For any character of the form(G15.1), its Weil term has
codimension four. Writing $Q(z)=q_0+q_1z+q_2z^2+\cdots$
and $r=q_0\ne0$, direct expansion gives
\[
       C\bigl(e^B(Q(\theta)+\lambda t)\bigr)
               =\left(q_2-\frac{q_1^2}{2r}\right)\theta^2.
                                                               \tag{G15.13}
\]
This formula also proves that $C$ is unchanged by
multiplication by any divisor exponential, so it applies
to $v_\ell$ with its displayed normalization.
The coefficient in(G15.13) cannot be zero here because
of the nonzero component(G15.12).

An ample real $(1,1)$ class on a complex torus of dimension
at least three is uniquely determined among ample classes
by its square. Indeed choose its positive invariant
Hermitian form as a basis form and diagonalize a second
positive invariant form relative to it. Equality of
squares says $\lambda_i\lambda_j=1$ for all $i\ne j$.
Positivity and three distinct indices force every
$\lambda_i=1$. Invariant forms represent the exterior
cohomology faithfully, proving the asserted uniqueness
as a cohomological statement.

If(G15.1) held, invariance of $C(v_\ell)$ under $\iota$
and(G15.13) would give
$(\iota^*\theta)^2=\theta^2$. Both classes are ample,
so the just-proved uniqueness gives $\iota^*\theta=\theta$.
In the degree-two K\"unneth decomposition, $\iota$
acts negatively precisely on
$H^1(X)\otimes H^1(\widehat X)$.
Thus $\theta$ has no mixed component and is a sum
of $(1,1)$ classes pulled from the two factors.
Its square has zero component in(G15.12), contradicting
the coefficient $1/8$ computed above.
This proves the claimed whole-character obstruction.

\paragraph{Certificate and scope.}
The independent script \texttt{point\_square\_character.py}
checks all four identities(G15.8) exactly over
$\mathbb Q(\ell)$, multiplies them to obtain the full
488-term character, and verifies(G15.11), inversion
invariance, the value at zero, and the holomorphic
coefficient $1/8$. Its JSON records the entire character.
The Ext, trace, ample-square and functorial arguments
are proved in the text, not inferred from the script.

The point ideal is simple, so this example does not
claim to supply a nonsimple indecomposable object.
It establishes an actual cone with all-degree trace
injectivity, identifies why the rank test leaves its
square open, and then excludes the specific equal-twist
Fourier construction(G15.5) by its full character.
It makes no assertion that every other Fourier
equivalence, unequal twist, extension, or CM
eightfold object has the same obstruction.
\hfill$\square$
